本プロジェクトの目標は，既約結晶的ルート系の分類定理をLean 4を用いて形式化することである：
\begin{theorem*}
  \label{thm:classification_root}
  \lean{classification_root}
  \leanok
  \uses{}
  $P$はcrystallographicなroot pairingであるとする．
  $C$を$P$から定まるCartan行列とする．
  このとき，$C$はCartan行列$A_n, B_n, C_n, D_n, E_6, E_7, E_8, F_4, G_2$のいずれか1つとCartan同値である．
\end{theorem*}
結晶的ルート系とは特定の条件を満たすユークリッド空間内のベクトルからなる集合であり，$A_n, B_n, \ldots, G_2$型に分類される．
数学の文脈では，Coxeter-Dynkin図形というグラフを定義し，その正定値性や部分グラフという概念を用いて分類できる．
具体的には，以下の性質を用いる：
\begin{enumerate}
  \item ルート系のCoxeter-Dynkin図形の正定値性
  \item 正定値グラフの部分グラフの正定値性
  \item $A_n, B_n, \ldots, G_2$の正定値性
  \item $\widetilde{A}_n, \widetilde{B}_n, \ldots, \widetilde{G}_2$の非正定値性
\end{enumerate}
これらを用いて，ルート系のCoxeter-Dynkin図形$\Gamma$を以下のように分類できる：
\begin{align}
  \Gamma \textrm{は正定値}
  &\implies \Gamma \textrm{の部分グラフは正定値} \\
  &\implies \Gamma \textrm{は} \widetilde{A}_n, \widetilde{B}_n, \ldots, \widetilde{G}_2 \textrm{を部分グラフに持たない} \\
  &\implies \Gamma \textrm{は} A_n, B_n, \ldots, G_2 \textrm{のいずれか}
\end{align}

しかし，本プロジェクトのLeanのコードでは，この流れで形式化していない．
なぜなら，部分グラフを用いた分類がうまくいかなかったからである．
上記性質の2, 3, 4は形式化できた．
1もできたと仮定して(証明の中身は\texttt{sorry}として)分類を試みたが，型の違いや添え字の付け方による同一視が難しかった．

そこで，Coxeter-Dynkin図形を登場させない方針に切り替えた．
代わりに，一般Cartan行列というものを定義し分類を行う．
一般Cartan行列からグラフを描くことができ，それはCoxeter-Dynkin図形とほぼ同じ情報を持っている．
逆に，Coxeter-Dynkin図形から一般Cartan行列を定めることもできる．
一般Cartan行列を用いた証明，すなわちLeanにおける証明は以下の流れで行われる：

正定値対称化一般Cartan行列から定まるグラフに関して，
\begin{enumerate}
  \item 木である(サイクルがない)ことを示す．
  \item 分岐は高々1つであることを示す．
  \item 各頂点は3分岐以下であることを示す．
  \item 3分岐のとき，格枝の長さを制限する．
\end{enumerate}

性質4の枝の長さについて，分岐点を含まない格枝の長さを$a, b, c$とすると，
\begin{equation}
  \frac{1}{a + 1} + \frac{1}{b + 1} + \frac{1}{c + 1} > 1
\end{equation}
という不等式が成り立つ．
Leanによる形式化でもこの不等式を証明している．
この不等式を導出するには，Schur complementを用いる方法が最も簡潔である．
その方針では$A_n^{-1}$の情報が必要であるが，これをLeanで求めるのは難しい．
そのため，Schur complementは用いない．
性質1, 2, 3, 4のいずれも，ベクトル$x$で，対称化一般Cartan行列$S(C)$に関する二次形式$\transpose{x} S(C) x$が非正になるものを構成することにより，矛盾を導くという方法を採っている．


なお，この文書では，数学の文脈に合わせて添え字を$1$始まりとしているが，Leanファイル内では$0$始まり，すなわち各添え字はこの文書から$-1$されていることに注意されたい．
つまり，Leanファイル内で$\texttt{Fin n} = \{0, 1, \ldots, n-1\}$と書かれている場合は，PDFファイルでは$I = \{1, 2, \ldots, n\}$と書いている．
